\NSPage{001}%
\begin{center}
  {\Large\bfseries FINITE TIME BLOWUP FOR NAVIER--STOKES\par}
  \vspace{2.5em}
  {\large OPENAI\par}
\end{center}

\vspace{2em}
\noindent\textsc{Abstract.} For every positive viscosity, we construct a solution of the three-dimensional
incompressible Navier--Stokes equations that starts from rest and develops unbounded velocity
in finite time while maintaining uniformly bounded kinetic energy.

\vspace{3em}
\begin{center}
  {\large\textsc{Contents}}
\end{center}
\begin{tabular*}{\textwidth}{@{\extracolsep{\fill}}llr}
1. & Introduction & 1\\
2. & Physical description of the blowup & 3\\
3. & Proof outline & 6\\
4. & Constructing the leading order flow & 24\\
5. & Correcting the base flow to every order & 45\\
6. & Auxiliary torus and separation of oscillatory supports & 62\\
7. & Oscillatory realization and correction of the residual stress & 73\\
8. & Compactly supported mean corrections & 88\\
9. & Residual improvement and the local field & 100\\
10. & Compact forcing and whole-space breakdown & 116\\
Appendix A. & Matching radial moments and constructing the heat exterior & 126\\
Appendix B. & Analytic profiles near the axis and their continuation & 144\\
Appendix C. & Realizing the admissible stress cone & 157\\
& References & 165
\end{tabular*}

\section{Introduction}\label{sec:1}

We consider the three-dimensional incompressible Navier--Stokes equations. A fundamental
question about the nature of these equations is whether a solution starting from
smooth data can develop unbounded velocity in finite time while its kinetic energy remains
bounded. We construct such a flow with zero initial velocity and a smooth force compactly
supported in space and time.

\begin{theorem}\label{thm:1.1}
For every \NSi{NS-F-P001-001}{\nu>0} there exist a force
\NSi{NS-F-P001-002}{f\in C_c^\infty(\mathbb R^3\times(0,\infty);\mathbb R^3)}, a compact set
\NSi{NS-F-P001-003}{K\subset\mathbb R^3}, and smooth velocity and pressure fields
\NSi{NS-F-P001-004}{u,p} on \NSi{NS-F-P001-005}{\mathbb R^3\times[0,1)} satisfying
\NSFormula{NS-F-P001-006}{display}{1.1}
\begin{equation}
  \partial_tu+(u\cdot\nabla)u-\nu\Delta u+\nabla p=f,
  \qquad \nabla\cdot u=0,
  \qquad u(\cdot,0)=0,
  \tag{1.1}\label{eq:1.1}
\end{equation}
such that \NSi{NS-F-P001-007}{\supp u(\cdot,t)\cup\supp p(\cdot,t)\subset K} for every
\NSi{NS-F-P001-008}{0\leq t<1},
\NSFormula{NS-F-P001-009}{display}{none}
\begin{equation*}
  \sup_{0\leq t<1}\lVert u(t)\rVert_{L^2(\mathbb R^3)}<\infty,
  \qquad
  \limsup_{t\uparrow1}\lVert u(t)\rVert_{L^\infty(\mathbb R^3)}=\infty.
\end{equation*}
Consequently, there is no smooth solution \NSi{NS-F-P001-010}{(u,P)} on
\NSi{NS-F-P001-011}{\mathbb R^3\times[0,\infty)} with the same force and initial
datum whose kinetic energy is uniformly bounded
\NSi{NS-F-P001-012}{\sup_{t\geq0}\frac12\int_{\mathbb R^3}|u(x,t)|^2\,dx<\infty}.
\end{theorem}

This establishes alternative (C) in the Millennium problem statement for Navier--Stokes as
stated by Fefferman in \cite{ref-13}. Compact support also yields the corresponding construction on
\NSi{NS-F-P001-013}{\mathbb T^3=\mathbb R^3/\mathbb Z^3}, establishing alternative (D) in \cite{ref-13}; see
Corollary~\ref{cor:10.6}.

\NSPage{002}%
\subsection{Historical context and previous work}\label{sec:1.1}
Incompressible and inviscid flows are described
by the classical Euler equations \cite{ref-12}:
\NSFormula{NS-F-P002-001}{display}{1.2}
\begin{equation}
  D_tu:=\partial_tu+(u\cdot\nabla)u=-\nabla p+f,
  \qquad \nabla\cdot u=0.
  \tag{1.2}\label{eq:1.2}
\end{equation}
The material acceleration \NSi{NS-F-P002-002}{D_tu} describes the acceleration of a parcel of fluid along its trajectory,
due both to the changing flow and the parcel's changing position. Crucially, the latter
contribution \NSi{NS-F-P002-003}{(u\cdot\nabla)u} is nonlinear in the velocity \NSi{NS-F-P002-004}{u}, creating most of the difficulty in the
analysis. The first equation in \eqref{eq:1.2} says that this acceleration is related by Newton's second
law to the sum of the negative pressure gradient and the external force. The second equation
tells us that the fluid is incompressible: we cannot squeeze the fluid in one direction without
simultaneously stretching it in another.

The Navier--Stokes equations \cite{ref-18,ref-21}, given in \eqref{eq:1.1}, add a viscous force that resists shear
between neighbouring fluid parcels. The viscous force is linear in \NSi{NS-F-P002-005}{u}. However, it significantly
changes the character of the flows by smoothing velocity gradients. Whether this smoothing
is always sufficient to prevent a finite-time singularity in three dimensions lies at the heart
of the regularity problem.

In 1934, Leray \cite{ref-16} constructed global finite-energy weak solutions of the three-dimensional
equations satisfying an energy inequality. Whether solutions starting from smooth data
remain smooth was left unresolved. Caffarelli, Kohn, and Nirenberg \cite{ref-5} proved that the
singular set of a suitable weak solution has zero one-dimensional parabolic Hausdorff
measure, a bound that allows isolated singularities. Escauriaza, Seregin, and Šverák \cite{ref-11}
proved regularity for the unforced Cauchy problem under a bounded scale-invariant
\NSi{NS-F-P002-006}{L_t^\infty L_x^3} norm.

Tao \cite{ref-22} constructed finite-time blowup for an averaged Navier--Stokes equation that
retains the energy cancellation of the original nonlinearity. For the original equations,
Buckmaster and Vicol \cite{ref-4} used convex integration to establish nonuniqueness among finite-
energy weak solutions. Albritton, Brué, and Colombo \cite{ref-1} subsequently constructed distinct
suitable Leray--Hopf solutions with zero initial velocity and the same force, using an unstable
vortex in similarity variables. Their force lies in \NSi{NS-F-P002-007}{L_t^1L_x^2} and is singular at the initial time.

The methods of Lifschitz and Hameiri and of Friedlander and Vishik \cite{ref-17,ref-14} describe
the evolution of wavevectors and velocity polarizations along a background flow. Craik
and Criminale \cite{ref-9} constructed exact finite-amplitude waves on affine background flows,
exploiting the cancellation of the wave's quadratic self-interaction. Further precedents for the
wave dynamics include the centrifugal-instability criteria of Leibovich and Stewartson \cite{ref-15}
and of Billant and Gallaire \cite{ref-2,ref-3} and the exact viscous shearing waves of Singh and Sridhar
\cite{ref-19}. The use of oscillations to realize a prescribed stress is central to the Euler constructions
of Daneri and Székelyhidi \cite{ref-10}.

Córdoba and Martínez-Zoroa \cite{ref-6} constructed finite-time singularities for the forced three-
dimensional Euler equations through successive amplification of increasingly concentrated
vortex layers. Córdoba, Martínez-Zoroa, and Zheng \cite{ref-8} extended this approach to hypodissipative
Navier--Stokes with small positive dissipation orders and forcing in a local
well-posedness class. These works established a strategy for singularity formation based on
amplification across scales while controlling the regularity of the external force. In their constructions,
larger-scale strain amplifies smaller-scale vorticity, with leading self-interactions
and feedback on the larger scales suppressed. For the incompressible porous media equation,
Córdoba and Martínez-Zoroa \cite{ref-7} constructed singularities from smooth initial data by successive
amplification of oscillatory layers, using approximations of increasing order to keep
every spatial derivative of the source uniformly bounded. Our construction also exploits
\NSPage{003}%
dynamical amplification, with a different role for the amplified disturbances: oscillatory
pulses generate a mean momentum flux that supplies the missing force on a collapsing
background vortex.

\section{Physical description of the blowup}\label{sec:2}

We first describe the underlying physics of the flow realizing Theorem~\ref{thm:1.1}; we then provide
a full sketch of the proof in Section~\ref{sec:3}.

For any incompressible flow \NSi{NS-F-P003-001}{u} and pressure \NSi{NS-F-P003-002}{p}, we can always define the external force
\NSi{NS-F-P003-003}{f} to be the residual in \eqref{eq:1.1}. The Navier--Stokes equations then hold by construction. The
challenge is to choose a flow that blows up while this residual remains smooth. The
individual terms in the momentum residual can diverge, but we must arrange sufficient
cancellation that their sum and all its derivatives extend smoothly through the singular time.

Specifically, we construct a vortex whose leading profile is self-similar, with radial width
decreasing faster than axial length. Its central flow combines inward spiralling with axial
outflow. As the core contracts, the azimuthal and axial velocity scales increase, strengthening
the shear in the surrounding annulus.

In the inner core, the profile equations impose the leading momentum balance. Joining
this flow to a smooth exterior whose speed decreases with radius leaves a nonzero residual
in the intervening annulus. For the background alone, the momentum residual becomes
unbounded as \NSi{NS-F-P003-004}{t\uparrow1}, so it cannot supply the smooth force required by the theorem. We add
spatially oscillatory pulses whose nonlinear momentum fluxes cancel the singular part of
this residual. Further corrections remove the remaining singular errors, leaving a force that
extends smoothly through \NSi{NS-F-P003-005}{t=1}.

\subsection{Inner core}\label{sec:2.1}
We construct our flow so that the singularity forms at the spatial origin at
time \NSi{NS-F-P003-006}{t=1}. We use cylindrical coordinates about the vertical axis:
\NSFormula{NS-F-P003-007}{display}{none}
\begin{equation*}
  x=(r\cos\theta,r\sin\theta,z),
\end{equation*}
where \NSi{NS-F-P003-008}{r} is distance from the axis, \NSi{NS-F-P003-009}{\theta} is angle around it, and \NSi{NS-F-P003-010}{z} is height. The velocity
components \NSi{NS-F-P003-011}{u_r}, \NSi{NS-F-P003-012}{u_\theta}, and \NSi{NS-F-P003-013}{u_z} describe motion toward or away from the axis, around the axis,
and along it, respectively. The leading flow \NSi{NS-F-P003-014}{u^{(0)}} is axisymmetric: its components depend on
\NSi{NS-F-P003-015}{r,z,t}, but not on \NSi{NS-F-P003-016}{\theta}.

In the central part of the inner region, fluid spirals inward toward the axis and flows
upward and downward on opposite sides of a dividing layer close to \NSi{NS-F-P003-017}{z=0} (Figure~\ref{fig:1}).
Farther above and below this central region, the inner profile also contains radial outflow.
Pressure decreases toward the axis, supplying the leading centripetal force. Its axial gradient
enters the axial momentum balance together with transport and viscosity. Near the axis, the
axial pressure force points toward the plane \NSi{NS-F-P003-018}{z=0}.

Radial inflow contributes to spin-up by carrying angular momentum toward smaller radii.
For a fluid particle experiencing no torque, the angular momentum per unit mass
\NSi{NS-F-P003-019}{ru_\theta} is conserved, so moving inward increases \NSi{NS-F-P003-020}{u_\theta}. In the actual core, viscosity transports angular
momentum outwards; the increasing characteristic speed reflects the balance between inward
transport and viscous loss. Incompressibility prevents the incoming fluid from accumulating
near the axis: the axial outflow carries it away, allowing the inward motion and spin-up to
continue.

To leading order, the core evolves self-similarly: its velocity profile has a fixed shape when
distances and velocities are measured in their respective time-dependent scales. The radial
\NSPage{004}%
\begin{figure}[t]
  \centering
  \NSFigurePlaceholder{NS-FIG-1}
  \caption{The central part of the inner core at successive times
  \NSi{NS-F-P004-001}{t_1<t_2<t_3<1}. In this region, fluid spirals inward and flows axially on opposite sides of a
  dividing layer close to \NSi{NS-F-P004-002}{z=0}. As the singular time approaches, the region of
  intense flow shrinks and the velocity increases. Its radius shrinks faster than
  its height; this difference is exaggerated in the schematic.}\label{fig:1}
\end{figure}%
and axial scales shrink at different rates. Writing \NSi{NS-F-P004-003}{\tau=1-t} for the time remaining before the
singularity, we have
\NSFormula{NS-F-P004-004}{display}{none}
\begin{equation*}
  \ell_r\asymp\tau^{1/2},
  \qquad \ell_z\asymp\tau^{1/2-h},
  \qquad 0<h<1/100,
\end{equation*}
where \NSi{NS-F-P004-005}{h} is fixed and \NSi{NS-F-P004-006}{\asymp} denotes bounds above and below by positive factors independent
of \NSi{NS-F-P004-007}{\tau}. Both lengths tend to zero, while
\NSi{NS-F-P004-008}{\ell_r/\ell_z\asymp\tau^h\to0}: the core becomes an increasingly
slender column concentrated at the origin, with volume of order
\NSi{NS-F-P004-009}{\tau^{3/2-h}}.

The characteristic velocity magnitudes of the leading flow satisfy
\NSFormula{NS-F-P004-010}{display}{none}
\begin{equation*}
  |u_\theta^{(0)}|,|u_z^{(0)}|\asymp\tau^{-1/2-h},
  \qquad |u_r^{(0)}|=O(\tau^{-1/2}).
\end{equation*}
These are typical scales across the core; in particular, the azimuthal velocity vanishes on
the axis. The flow is smooth for every \NSi{NS-F-P004-011}{t<1}, but its largest speeds become unbounded in
a shrinking neighbourhood of the origin as \NSi{NS-F-P004-012}{t\uparrow1}. The total kinetic energy of the core is of
order \NSi{NS-F-P004-013}{\tau^{1/2-3h}}, which tends to zero despite the increasing speeds.

The different velocity scales give different Reynolds numbers. For a speed
\NSi{NS-F-P004-014}{U}, length \NSi{NS-F-P004-015}{\ell},
and viscosity \NSi{NS-F-P004-016}{\nu}, the Reynolds number \NSi{NS-F-P004-017}{U\ell/\nu} compares the viscous diffusion time
\NSi{NS-F-P004-018}{\ell^2/\nu} with
the time \NSi{NS-F-P004-019}{\ell/U} for fluid motion over that distance. Using the core radius as the length scale,
the characteristic angular and radial Reynolds numbers are
\NSFormula{NS-F-P004-020}{display}{none}
\begin{equation*}
  \mathrm{Re}_\theta:=\frac{|u_\theta|\ell_r}{\nu}\asymp\tau^{-h}\longrightarrow\infty,
  \qquad
  \mathrm{Re}_r:=\frac{|u_r|\ell_r}{\nu}=O(1),
\end{equation*}
for fixed \NSi{NS-F-P004-021}{\nu>0}.

The angular Reynolds number grows without bound: fluid makes increasingly many
turns during a radial diffusion time. The radial Reynolds number remains bounded, so
viscosity continues to compete with radial inflow. The corresponding transport and diffusion
\NSPage{005}%
\begin{figure}[t]
  \centering
  \NSFigurePlaceholder{NS-FIG-2}
  \caption{Pulse geometry and perturbation velocities. \emph{Left:} At fixed height,
  pulse envelopes occupy complete rings; the close-ups show radial and azimuthal
  perturbation velocities. \emph{Right:} A radial--axial section, where each
  localized patch is a cross-section of a complete ring; the close-ups show radial
  and axial perturbation velocities. Colors distinguish the two families. The
  specific pulse arrangement shown is schematic; active locations, spacing, and
  aspect ratios are illustrative.}\label{fig:2}
\end{figure}%
rates satisfy
\NSFormula{NS-F-P005-001}{display}{none}
\begin{equation*}
  \frac{|u_r|}{\ell_r}=O(\tau^{-1}),
  \qquad \frac{|u_z|}{\ell_z}\asymp\tau^{-1},
  \qquad \frac{\nu}{\ell_r^2}\asymp\tau^{-1}.
\end{equation*}
Thus radial diffusion and transport by the radial and axial flow remain in the leading balance
as the core contracts. Axial diffusion is weaker: the ratio of axial to radial diffusion rates is
\NSFormula{NS-F-P005-002}{display}{none}
\begin{equation*}
  \frac{\nu/\ell_z^2}{\nu/\ell_r^2}
  =\frac{\ell_r^2}{\ell_z^2}
  \asymp\tau^{2h}\longrightarrow0.
\end{equation*}

The core must also provide enough shear at its outer edge to amplify the pulses placed
there. Away from the middle plane, axial flow carries angular momentum from more rapidly
rotating layers, supporting a steep radial decrease in angular velocity. With exact reflection
symmetry, this transport would vanish at \NSi{NS-F-P005-003}{z=0}. Symmetry would also impose
\NSi{NS-F-P005-004}{u_z(r,0,t)=0},
leaving no radial shear of the axial velocity to compensate. We therefore choose a slightly
asymmetric axial profile, with a small upward bias and nonzero velocity at
\NSi{NS-F-P005-005}{z=0}. This
separates the layer where rotational amplification is weak from the layer where axial shear
vanishes. The radial shear of axial velocity then supplies the additional amplification needed
near the middle plane; farther away, the rotational mechanism already suffices.

\subsection{Annulus and pulses}\label{sec:2.2}
At the edge of the core, there is an annular region where the inner
vortex joins onto an exterior flow. The background satisfies its leading momentum equations
inside the core, but this transition leaves an imbalance in the annulus. The external force
needed to sustain the background becomes unbounded as \NSi{NS-F-P005-006}{t\uparrow1}, so it cannot serve as the
smooth external force required by the construction.

We therefore perturb the velocity so that the fluid's own motion supplies the missing
momentum transport. Through the nonlinear advection term, the perturbation produces
additional momentum fluxes whose spatial variation acts as an internal force on the background.
We require this contribution to cancel the singular part of the background residual,
leaving only a remainder that extends smoothly through \NSi{NS-F-P005-007}{t=1}. These internal forces
redistribute momentum within the fluid.

\NSPage{006}%
Our perturbation consists of a sequence of spatially oscillatory pulses that are localized
in radius, height, and time, but extend around a complete ring. An exponentially small
external force seeds each pulse; the background shear supplies its subsequent growth. Their
cylindrical velocity components have zero angular average, but their momentum fluxes
need not. For example, outward motion carrying an azimuthal velocity surplus and inward
motion carrying a deficit both increase outward angular momentum flux. Reversing both
velocity components leaves their product unchanged. Axial momentum is transported in the
same way. We arrange the net transport into or out of each region to supply the contribution
missing from the background's momentum balance.

The pulses grow by extracting energy from the background shear. A simple example
illustrates how this amplification occurs. In a rotating flow, an azimuthal velocity surplus
produces outward acceleration. If angular velocity decreases sufficiently rapidly with radius,
this outward motion increases the surplus relative to the surrounding fluid. The two effects
reinforce one another, giving exponential growth when the amplification exceeds viscous
damping. In the full flow, axial shear also contributes to the amplification.

The same shear also changes the oscillation's spatial pattern: fluid at different radii carries
it along at different speeds. We choose the pulse orientations so that this progressively
shortens the radial wavelength. The changing orientation weakens the amplification, while
shortening the radial wavelength increases the total squared wavenumber and hence viscous
damping. We choose the initial wavelength so that amplification dominates at first, but
damping overtakes it later. The tails are sufficiently small that the cutoff errors and all their
derivatives vanish at the singularity.

We combine two families of pulses with different ratios of radial angular-momentum
flux to radial axial-momentum flux. The background profile and oscillation directions are
chosen together so that the averaged quadratic velocity products of the two families supply
both components of the required leading stress. We arrange the pulses on successively finer
spatial and temporal scales as the singularity approaches. Further corrections cancel the
remaining singular terms in the momentum residual, giving a flow whose velocity becomes
unbounded while its residual force extends smoothly through the singular time.

\subsection{Exterior flow and localization}\label{sec:2.3}
Beyond the annulus, the flow is purely azimuthal and
independent of height: fluid circles the axis, with speed decreasing as the radius increases.
Pressure supplies the centripetal acceleration, while the azimuthal velocity evolves by viscous
diffusion. We choose this outer flow to satisfy the radial heat equation for the azimuthal
velocity exactly, and call it the heat exterior. Its momentum residual vanishes, so no external
force is needed to sustain this part of the local flow.

At every fixed positive radius, the exterior velocity and pressure have smooth limits,
together with all their derivatives, as the singular time approaches. Together with the
regularity of the inner construction away from the origin, this allows us to cut off the flow
smoothly in space. The additional force produced by this cutoff remains smooth through the
singular time, while the cutoff leaves the concentrating core unchanged.

\section{Proof outline}\label{sec:3}

The physical description in Section~\ref{sec:2} leads to a precise cancellation problem. We must
construct a velocity and pressure whose momentum residual extends smoothly through the
singular time, while the velocity becomes unbounded. At viscosity one, the residual is
\NSFormula{NS-F-P006-001}{display}{3.1}
\begin{equation}
  \mathscr R(u,p)=\partial_tu+(u\cdot\nabla)u-\Delta u+\nabla p.
  \tag{3.1}\label{eq:3.1}
\end{equation}
